% from ega3-1-fr.tex
\oldpage[III]{86}
\item $A=\Gamma(X,\mathcal O_X)$, $M=\Gamma(X,\mathcal F)$~;
\item $\mathbf f=(f_i)_{1\leq i\leq r}$ un système fini d'éléments de $A$~;
\item $U_i=X_{f_i}$ l'ensemble ouvert
(\hyperref[0.5.5.2-fr]{0\textsubscript{I}, 5.5.2}) des $x\in X$ tels que
$f_i(x)\neq0$~;
\item $U=\bigcup_{i=1}^r U_i$~;
\item $\mathfrak U$ le recouvrement $(U_i)_{1\leq i\leq r}$ de $U$.
\end{itemize}
\end{notation}

\begin{env}[1.2.2]
\label{III.1.2.2-fr}
Supposons que $X$ soit, ou bien un préschéma dont l'espace sous-jacent est
\emph{noethérien}, ou bien un \emph{schéma} dont l'espace sous-jacent soit
\emph{quasi-compact}. On sait alors (\hyperref[I.9.3.3-fr]{I, 9.3.3}) que
l'on a $\Gamma(U_i,\mathcal F)=M_{f_i}$. Nous poserons
\[
U_{i_0i_1\ldots i_p}=\bigcap_{k=0}^p U_{i_k}
=X_{f_{i_0}f_{i_1}\cdots f_{i_p}}
\]
(\hyperref[0.5.5.3-fr]{0\textsubscript{I}, 5.5.3})~; on a donc aussi
\begin{equation}
\label{III.1.2.2.1-fr}
\tag{1.2.2.1}
\Gamma(U_{i_0i_1\ldots i_p},\mathcal F)=M_{f_{i_0}f_{i_1}\cdots f_{i_p}}.
\end{equation}
Or (\hyperref[0.1.6.1-fr]{0\textsubscript{I}, 1.6.1})
$M_{f_{i_0}f_{i_1}\cdots f_{i_p}}$ s'identifie à la limite inductive
$\varinjlim_n M_{i_0i_1\ldots i_p}^{(n)}$, où le système inductif est formé
des $M_{i_0i_1\ldots i_p}^{(n)}=M$, l'homomorphisme
$\varphi_{nm}:M_{i_0i_1\ldots i_p}^{(m)}\to
M_{i_0i_1\ldots i_p}^{(n)}$ étant la multiplication par
$(f_{i_0}f_{i_1}\cdots f_{i_p})^{n-m}$ pour $m\leq n$. Désignons par
$C_n^p(M)$ l'ensemble des applications \emph{alternées} de $[1,r]^{p+1}$
dans $M$ (pour tout $n$)~; ces $A$-modules forment encore un système inductif
pour les $\varphi_{nm}$. Si $C^p(\mathfrak U,\mathcal F)$ est le groupe des
$p$-cochaînes \emph{alternées} de Čech relatif au recouvrement $\mathfrak U$,
à coefficients dans $\mathcal F$ (G, II, 5.1), il résulte de ce qui précède
que l'on peut écrire
\begin{equation}
\label{III.1.2.2.2-fr}
\tag{1.2.2.2}
C^p(\mathfrak U,\mathcal F)=\varinjlim_n C_n^p(M).
\end{equation}
Or, avec les notations de (\hyperref[III.1.1.2-fr]{1.1.2}), $C_n^p(M)$
s'identifie à $K^{p+1}(\mathbf f^n,M)$, et l'application $\varphi_{nm}$
s'identifie à l'application $\varphi_{\mathbf f^{n-m}}$ définie dans
(\hyperref[III.1.1.6-fr]{1.1.6}). On a donc, pour tout $p\geq0$, un
isomorphisme canonique fonctoriel en $\mathcal F$
\begin{equation}
\label{III.1.2.2.3-fr}
\tag{1.2.2.3}
C^p(\mathfrak U,\mathcal F)\xrightarrow{\sim}C^{p+1}((\mathbf f),M).
\end{equation}
En outre, la formule (\hyperref[III.1.1.2.3-fr]{1.1.2.3}) et la définition de
la cohomologie d'un recouvrement (G, II, 5.1) montrent que les isomorphismes
(\hyperref[III.1.2.2.3-fr]{1.2.2.3}) sont compatibles avec les opérateurs
cobords.
\end{env}

\begin{proposition}[1.2.3]
\label{III.1.2.3-fr}
Si $X$ est un préschéma dont l'espace sous-jacent est noethérien, ou un
schéma dont l'espace sous-jacent est quasi-compact, il existe un isomorphisme
canonique fonctoriel en $\mathcal F$
\begin{equation}
\label{III.1.2.3.1-fr}
\tag{1.2.3.1}
H^p(\mathfrak U,\mathcal F)\xrightarrow{\sim}H^{p+1}((\mathbf f),M)
\qquad\text{pour tout }p\geq1.
\end{equation}
On a en outre une suite exacte fonctorielle en $\mathcal F$
\begin{equation}
\label{III.1.2.3.2-fr}
\tag{1.2.3.2}
0\longrightarrow H^0((\mathbf f),M)\longrightarrow M
\longrightarrow H^0(\mathfrak U,\mathcal F)
\longrightarrow H^1((\mathbf f),M)\longrightarrow0.
\end{equation}
\end{proposition}

